With Corollary \ref{Core} as motivation, define for any $\ii,\jj \in \mathcal{W}_m$,
\[
\tilde{e}_{\jj \ii} = \sum_{\tau \in D_{\jj}} q^{-d_{\ii}(\tau)}e_{\tau(\jj)\tau(\ii)} .
\]

\subsection{Proof of Theorem \ref{Main}}
We will now prove Theorem \ref{Main}. Begin by rewriting $R_{0,\{1,\dots,m\}}$. By definition,
\begin{align*}
R_{0,\{1,\dots,m\}} &= R_{0m}R_{0m-1}\cdots R_{01} \\
&=\bigg( \sum_{i_m \geq j_m} \hat{E}_{i_mj_m} \otimes \id^{\otimes m-1} \otimes e_{j_mi_m} \bigg) \cdots \bigg(\sum_{i_1\geq j_1}\hat{E}_{i_1j_1}\otimes e_{j_1i_1} \otimes \id^{\otimes m-1}\bigg)
\\&= \sum_{i_m \geq j_m} \cdots \sum_{i_1\geq j_1} \hat{E}_{i_mj_m} \cdots \hat{E}_{i_1j_1} \otimes e_{j_1i_1} \otimes \dots \otimes e_{j_mi_m} \\
&= \sum_{\ii,\jj \in \mathcal{W}_m } \sum_{\zeta \in D_{\ii},\tau \in D_{\jj}}\hat{E}_{\bar{\zeta}(\ii)\bar{\tau}(\jj)}^{+} \otimes e_{\tau(\jj)\zeta(\ii)}.
\end{align*}
From here, the goal is to write this expression as an element of $\Uqgl \otimes \End\big(P_m\big(V^{\otimes m}\big)\big)$. By Corollary \ref{Core},
\[
R_{0,\{1,\dots,m\}} = \sum_{\ii,\jj \in \mathcal{W}_m } \sum_{\zeta \in D_{\ii},\tau \in D_{\jj}}q^{\inv(\sigma_{\ii}(\tau)) -\inv(\zeta)}\hat{E}^+_{\bar{\zeta}(\ii)\bar{\tau}(\jj)} \otimes e_{\tau(\jj)\sigma_{\ii}(\tau)(\ii)}.
\]
By definition, $\tau = \sigma_{\ii}(\tau) \xi_{\ii}(\tau)$ and $d_{\ii}(\tau) = \inv(\xi_{\ii}(\tau))$, so therefore
\[
R_{0,\{1,\dots,m\}} = \sum_{\ii,\jj \in \mathcal{W}_m } \sum_{\zeta \in D_{\ii},\tau \in D_{\jj}}q^{\inv(\tau) -d_{\ii}(\tau) -\inv(\zeta)}\hat{E}^+_{\bar{\zeta}(\ii)\bar{\tau}(\jj)} \otimes e_{\tau(\jj)\tau(\ii)}.
\]
By Lemma \ref{prop 1} and the identity $\inv(\bar{\tau}) = m(m-1)/2 - \inv(\tau)$,
\[
\sum_{\zeta \in D_{\ii}}q^{\inv(\tau) -\inv(\zeta)}\hat{E}_{\bar{\zeta}(\ii)\bar{\tau}(\jj)}
\]
does not depend on $\tau$. Therefore, the $R$-matrix factors as
\begin{align*}
R_{0,\{1,\dots,m\} } &= \sum_{\ii,\jj \in \mathcal{W}_m } \bigg( \sum_{\zeta \in D_{\ii}}q^{\inv(\tau) -\inv(\zeta)}\hat{E}^+_{\bar{\zeta}(\ii)\bar{\tau}(\jj)} \bigg) \otimes \bigg( \sum_{\tau \in D_{\jj}} q^{-d_{\ii}(\tau)}e_{\tau(\jj)\tau(\ii)} \bigg) \\
&= \sum_{\ii,\jj \in \mathcal{W}_m } \tilde{E}_{\ii\jj}^+ \otimes e_{\jj\ii}.
\end{align*}
A similar argument shows that
\[
R^T_{0,\{1,\dots,m\}}=\sum_{\ii,\jj \in \mathcal{W}_m } \tilde{E}^-_{\jj\ii} \otimes e_{\ii\jj}.
\]

Therefore,
\[
\Gamma_m = R^T_{0,\{1,\dots,m\} } R_{0,\{1,\dots,m\} } = \sum_{\ii,\jj \in \mathcal{W}_m } \sum_{\ii',\jj' \in \mathcal{W}_m } \tilde{E}_{\jj \ii}^-\tilde{E}_{\ii'\jj'}^+ \otimes e_{\ii\jj}e_{\jj'\ii'}.
\]
By Corollary \ref{Core},
\[
C_m = (\id \otimes \tr_q)(\Gamma_m) = \sum_{\ii,\jj \in \mathcal{W}_m } \tilde{E}_{\jj \ii}^-\tilde{E}_{\ii\jj}^+ \otimes \tr_q(e_{\ii\jj}e_{\jj\ii}).
\]

It just remains to calculate the quantum trace. It is given by
\[
\tr_q(e_{\ii\jj}e_{\jj\ii}) = \tr_q(e_{\ii\ii}) = \tr\big(q^{-NE_{00} -(N-2)E_{11} + \dots + NE_{NN}} e_{\ii\ii}\big) = q^{-N\mu_{0} -(N-2)\mu_{1} + \dots + N\mu_{N}},
\]
where $\mu=\mu^{(N)}(\ii)$. Applying Lemma \ref{Leem} finishes the proof.

